
\(F^p\wedge^n(\underline{E})=\wedge^n(\underline{E})\), et si \(p<0\), \(F^p\wedge^n(\underline{E})=0\). Ceci définit bien une filtration finie sur \(\wedge^n(\underline{E})\).

On définit un homomorphisme~:
\[
  f_p:\wedge^{n-p}(\underline{E}')\otimes_A\wedge^p(\underline{E}'')
  \longrightarrow F^p\wedge^n(\underline{E})/F^{p-1}\wedge^n(\underline{E})
\]
en posant \(f_p(x'_1\wedge\ldots\wedge x'_{n-p}\otimes x''_1\wedge\ldots\wedge x''_p)=\operatorname{cl}_{\operatorname{Gr}^p}\bigl(u(x'_1)\wedge\ldots\wedge u(x'_{n-p})\wedge x_1\wedge\ldots\wedge x_p\bigr)\) qui ne dépend visiblement pas du choix des \(x_i\) relevant les sections données \(x''_i\) de \(\underline{E}''\). Si la suite exacte est localement scindée, \(f_p\) est un isomorphisme, car localement \(\wedge^n(\underline{E})\simeq\bigoplus_{p+q=n}\wedge^p(\underline{E}')\otimes_A\wedge^q(\underline{E}'')\).

Dans le cas où \(\underline{E}',\underline{E},\underline{E}''\) sont localement facteurs directs de Modules libres de type fini, les \(\operatorname{Gr}^p\wedge^n(\underline{E})\) sont localement facteurs directs de libres de type fini, donc les \(F^p\wedge^n(\underline{E})\) également~; on peut donc écrire dans \(K\)~:
\(\lambda^n(E)=\sum_{p=0}^{n}\operatorname{cl}(\operatorname{Gr}^p\wedge^n(\underline{E}))=\sum_{p=0}^{n}\lambda^{n-p}(E')\mathbin{.}\lambda^p(E'')\), d'où l'additivité.

Soient \((T,\underline{A})\), \((T',\underline{A}')\) deux topos commutativement annelés, et \(f:(T,\underline{A})\longrightarrow(T',\underline{A}')\) un morphisme de topos annelés. On sait (Exp. IV, 2.7) que \(f\) définit un morphisme d'anneau sur les \(K\) correspondants, soit \(f^*:K^{\bullet}_{\naif}(T',\underline{A}')\longrightarrow K^{\bullet}_{\naif}(T,\underline{A})\). Comme l'image inverse commute aux puissances extérieures, le morphisme \(f^*\) correspondant sur les \(K\) est un \(\lambda\)-homomorphisme pour la structure qui vient d'être définie.

Signalons deux cas particuliers :
\begin{itemize}[label=-]
\item Soit \(X\) un schéma~; \(K^{\bullet}_{\naif}(X)\) est ainsi muni d'une structure de pré-\(\lambda\)-anneau; c'est ce cas qui servira dans la suite du séminaire, et il peut être considéré comme « modèle » pour la théorie générale des 
